\documentstyle[%


righttag%


,11pt%


,amssymb%


,russian%               remove in Eng.


,izvru%                 Set izven% in Eng.


]{amsart}





%





\newcommand{\Diff}{\operatorname{Diff}}


\newcommand{\rank}{\operatorname{rank}}


\newcommand{\Grass}{\operatorname{Grass}}


\newcommand{\Gl}{\operatorname{Gl}}


\newcommand{\GL}{\operatorname{GL}}


\newcommand{\Hom}{\operatorname{Hom}}


\newcommand{\Ann}{\operatorname{Ann}}


\newcommand{\Ker}{\operatorname{Ker}}


\newcommand{\loc}{\operatorname{loc}}


\newcommand{\dom}{\operatorname{dom}}


\newcommand{\id}{\operatorname{id}}


\newcommand{\la}{\langle}


\newcommand{\ra}{\rangle}


\newcommand{\tts}[1]{}


\newtheorem{theorema}{Теорема}


\newtheorem{theoremb}{Теорема}


\newtheorem{theoremc}{Теорема}


\newcommand{\A}{{\cal A}}


\renewcommand{\a}{\alpha}


\renewcommand{\b}{\beta}


\renewcommand\c{\circ}


\newcommand{\ccc}{\c c}


\newcommand{\C}{{\cal C}}


\newcommand{\CC}{{\Bbb C}}


\newcommand{\dd}{\dotsc}


\renewcommand{\d}{\delta}


\newcommand{\D}{{\cal D}}


\newcommand{\e}{\eta}


\newcommand{\E}{{\cal E}}


\newcommand{\F}{{\cal F}}\newcommand{\Fg}{{\frak F}}


\newcommand{\ga}{\gamma}


\newcommand{\g}{{\frak g}}


\newcommand{\G}{{\cal G}}


\newcommand{\op}[1]{\mathop{\mathrm{#1}}\nolimits}


\renewcommand{\Gl}{\operatorname{Gl}}


\newcommand{\h}{{\frak h}}


\newcommand{\J}{{\cal J}}


\renewcommand{\l}{{\cal L}}


\newcommand{\Li}{\operatormane{Lie}}


\newcommand{\li}{\operatorname{lie}}


\renewcommand{\ll}{\lambda}


\newcommand{\La}{\Lambda}


\newcommand{\m}{\mu}


\renewcommand{\O}{\Omega}


\newcommand{\oo}{\omega}


\newcommand{\ok}{{\cal O}}


\newcommand{\ot}{\otimes}


\newcommand{\p}{\partial}


\renewcommand{\P}{\Phi}


\newcommand{\PP}{\Pi_N}


\newcommand{\R}{{\Bbb R}}


\newcommand{\St}{\op{St}^l_a}


\renewcommand{\t}{\tau}


\newcommand{\tf}{{\frak t}}


\newcommand{\te}{\theta}


\newcommand\ti{\widetilde}


\newcommand{\T}{\Theta}


\newcommand{\ups}{\upsilon}


\renewcommand{\u}{{\cal U}}


\newcommand{\ve}{\varepsilon}


\newcommand{\vg}{{\frak v}}


\newcommand{\vp}{\varphi}


\newcommand{\V}{{\cal V}}


\newcommand{\w}{{\cal W}}


\newcommand{\we}{\wedge}


\newcommand{\x}{\xi}


\newcommand{\Y}{{\cal Y}}


\newcommand{\z}{\sigma}


\newcommand{\Z}{{\Bbb Z}}





\theoremstyle{plain}


\theorembodyfont{\it}


\newtheorem{thmnonum}{\hskip\parindent Теорема}


\renewcommand{\thethmnonum}{}





\theoremstyle{definition}


\theorembodyfont{\rm}


\newtheorem{cors}{\hskip\parindent Следствие}[section]








%\hoffset=-15mm%                  For Eng.


\setcounter{page}{30}


\god{2004}


\nomer{11~(510)} %                        Remove in Eng.


%\nomer{Vol.\,48, No.\,11}%               Set in Eng.


%\str{\pageref{begin}--\pageref{end}}%  Set in Eng.


\udk{514.763}





\author{\it Б.С.\,Кругликов, В.В.\,Лычагин}


% NB:   ^^ remove in Eng.


\title{Транзитивные и  трансверсальные действия псевдогрупп на


подмногообразиях}


\thanks{}





\date{}


%\pagestyle{myheadings}%         Set in Eng.


\pagestyle{plain} %              Remove in Eng.


\begin{document}


%\thispagestyle{plain}%          Set in Eng.


\maketitle


\label{begin}











\section*{Введение}\label{sec0}








В данной работе рассматривается проблема эквивалентности


подмногообразий относительно транзитивного действия псевдогруппы.


Определяются соответствующие формальные инварианты ($l$-варианты и


$l$-коварианты). Исследуются их свойства,  и они вычисляются


в ряде важных случаев. Основное приложение о действии псевдогруппы Ли


преобразований в пространстве струй соответствует


проблеме эквивалентности дифференциальных уравнений.








Группы преобразований или более широко~--- псевдогруппы, изучение


которых было начато в \cite{Lie}, \cite{C},


играют центральную роль в геометрии и анализе.





Псевдогруппа $G \subset \Diff_{\loc}(M)$, действующая на


многообразии $M$, есть множество локальных диффеоморфизмов $\varphi$ с


областями определения $\dom(\varphi)$, удовлетворяющих следующим


свойствам:





1. eсли локальные диффеоморфизмы $\varphi$, $\psi$ принадлежат $G$,


то их композиция $\varphi \circ \psi$ со своей областью определения


тоже принадлежит $G$,





2. eсли $\varphi \in G$, то $\varphi^{-1} \in G$,





3. $\id_M \in G$,





4. $\varphi \in G$ тогда и только тогда, когда для любого открытого


подмножества $U \subset \dom (\varphi)$ ограничение


$\varphi|_U$ принадлежит $G$.





Говорят, что псевдогруппа $G$ имеет порядок $l$, если $l$ есть


наименьшее натуральное число со свойством, что локальный диффеоморфизм


$\varphi$ лежит в $G$ тогда и только тогда, когда для любой точки


$a \in \dom (\varphi)$ его $l$-струя $[\varphi]^l_a$ лежит в


подмногообразии $G^l$ пространства $l$-струй локальных


диффеоморфизмов многообразия $M$,


т.\,е. псевдогруппа определяется дифференциальным уравнением порядка


$l$.





Локальный диффеоморфизм $\varphi \in G$


определяет отображение ($l$-продолжение) $\varphi_{(l)} : J^l_r(M) \to


J^l_r(M)$


пространства $l$-струй подмногообразий $M$ коразмерности $r$,


которое удовлетворяет следующим свойствам:





\begin{itemize}


\item[---] $(\varphi \circ \psi)_{(l)} = \varphi_{(l)} \circ \psi_{(l)}$,


\item[---] $\id_{(l)} = \id^l_M$,


\item[---] $(\varphi^{-1})_{(l)}=\varphi_{(l)}^{-1}$


\end{itemize}


Эти свойства являются основой формального аспекта теории псевдогрупп


(см. \cite{E}--\cite{SS}).





В данной работе вводится более общее понятие инфинитезимальной


псевдогруппы. Именно, $l$-псевдогруппа есть псевдогруппа


преобразований пространства струй конечного порядка, причем


интегрируемость этих преобразований не предполагается.


Тем не менее, такая псевдогруппа конечного


порядка полезна для получения инвариантов дифференциальных уравнений и


кривизн геометрических структур.





Критерий формальной интегрируемости для инфинитезимальных псевдогрупп


основывается на алгебраической технике, описанной в данной


работе. Переход от формальной интегрируемости к локальной не может


быть произведен автоматически, и, вообще говоря, невозможен. Тем не


менее формальная интегрируемость влечет локальную в следующих случаях.





\begin{itemize}


\item[---]


Псевдогруппы конечного типа (символ $\g^k \equiv 0$ для больших


$k$). При этом условии проинтегрированная псевдогруппа конечномерна.


\item[---]


Аналитические псевдогруппы. Это следствие теоремы Картана--Келера,


которая выполняется для дифференциальных уравнений достаточно общего


вида \cite{M}, \cite{M1} \cite{KLV}.


\item[---]


Эллиптические псевдогруппы аналитического типа


(см. \cite{M}, \cite{M1}, \cite{S}).


\item[---]


Транзитивные плоские псевдогруппы (см. \cite{BM}, \cite{P}).


\end{itemize}





Проблема глобальной интегрируемости (или эквивалентности) может быть


решена только в очень специальных случаях


(см. \cite{S}, \cite{GS}, \cite{T}).








Цель данной статьи состоит в развитии теории инвариантов псевдогрупповых


действий на подмногообразиях. Для групп Ли это теория 


дифференциальных инвариантов. На уровне струй конечного порядка


получаются объекты, названные здесь $l$-ковариантами. Вычисление


$l$-ковариантов осуществляется с помощью некоторых когомологий


наподобие формальных когомологий Спенсера, устанавливается их связь


с проблемой эквивалентности подмногообразий относительно


псевдогруппового действия.


Псевдогруппы Ли состоят из псевдоавтоморфизмов геометрических


структур. Для таких псевдогрупп приводится ряд вычислений. Наиболее


важное приложение касается псевдогруппы преобразований Ли пространства


струй, т.\,е. преобразований, сохраняющих распределение Картана. Они


задают эквивалентность дифференциальных уравнений, поэтому предложенный


способ дает возможность получать инварианты дифференциальных уравнений.











\section{Формальные псевдогруппы}\label{sec1}





Пусть $M$~--- гладкое многообразие, обозначим через $J^l_r(M)$


соответствующее пространство струй. Точками этого пространства


являются $l$-струи $a_l = [N]_a^l$ подмногообразий $N \subset M$


коразмерности $r$ в точках $a \in M$. Пусть $\rho_{i,j} : J^i_r(M) \to


J^j_r(M)$, $i>j$,~--- естественные проекции. Слои этих проекций


наделены канонической аффинной структурой \cite{KLV}, \cite{Ly},


ассоциированной со структурой векторного пространства, описываемой


следующим образом (достаточно это сделать для


$\frak{F}(a_{l-1})=\rho^{-1}_{l,l-1}(a_{l-1})$).


Положим $\tf_a=T_aN=[N]^1_a$ и $\vg_a=T_aM/T_aN$.


Пусть $a_l\in J^l_r(M)$, $a_{l-1}=\rho_{l,l-1}(a_l)$.


Тогда


$T_{a_l}\Fg(a_{l-1})\simeq S^l\tf^*_a\ot\vg_a$,


и получаем точную последовательность


$$


0\to S^l\tf^*_a\ot\vg_a\to


T_{a_l}J^l_r(M)@>{(\rho_{l,l-1})_*}>>


T_{a_{l-1}}J^{l-1}_r(M)\to 0.


$$








Обозначим через $D^l(M)\subset J^l(M,M)$ открытое плотное


подмножество, состоящее из $l$-струй локальных диффеоморфизмов.


Если наделить его частично определенной операцией композиции, получим


пример $l$-псевдогруппы.





Для того чтобы определить понятие $l$-псевдогруппы в общем случае,


напомним понятие продолжения дифференциального уравнения из


геометрической теории дифференциальных уравнений


(подробности см. в \cite{KLV}, \cite{Ly}, \cite{Gu}).


Продолжение дифференциального уравнения


$\E\subset J^l_r(M)$ есть


$$


\E^{(1)}=\{a_{l+1}=[s]_a^{l+1}\mid \text{ струя } j_l(s) \text{ сечения


$s$ касательна к }\E\text{ в }a_l\}\subset


J^{l+1}_r(M).


$$


Эквивалентным образом это можно записать в виде


$\E^{(1)}=\{a_{l+1}\,|\,L(a_{l+1})\subset T_{a_l}\E\}$, где для


$a_{l+1}=[s]_a^{l+1}$ полагаем $L(a_{l+1})=T_{a_l}j_l(s)$,


$a_l=\rho_{l+1,l}(a_{l+1})$.








\begin{defn}\label{pseudogr}


$l$-псевдогруппа есть набор подрасслоений


$G^j\subset D^j(M)$, $0<j\le l$ такой, что


\begin{itemize}


\item[1.] eсли $\vp_j,\psi_j\in G^j$, то $\vp_j\circ\psi_j\in G^j$ там


где эта композиция определена,


\item[2.] $\id^j_M\in G^j$,


\item[3.] если $\vp_j\in G^j$, то $\vp_j^{-1}\in G^j$,


\item[4.] отображение $\rho_{j,j-1}:G^j\to G^{j-1}$ есть расслоение.


\end{itemize}


Также, по аналогии с геометрической теорией дифференциальных


уравнений, полагаем $G^0=J^0(M,M) = M\times M$, что эквивалентно


{\em транзитивности} псевдогруппового действия.





$l$-псевдогруппа называется $l$-интегрируемой если


$G^j\subset(G^{j-1})^{(1)}$ для всех $0<j\le l$.


\end{defn}





Интегрируемость определенных таким образом псевдогрупп


$G=\{G^j\}_{j=1}^l$ можно


изучать с помощью стандартного метода продолжений и проекций


(\cite{GS}, \cite{KLV}, \cite{S}).





Пусть $G^j_{a,b}=\{\vp_j\in G^j\mid\vp_0(a)=b\}$,


$G^j_a=G^j_{a,a}$ --- подгруппа в $G^j$ и ее нормальная подгруппа


$\frak{G}^j_a=\Ker[\rho_{j,j-1}:G^j_a\to G^{j-1}_a]$ 


является коммутативной при $j>1$, а при


$j=1$ имеем $\frak{G}^1_a=G^1_a\subset\Gl(T_aM)$.








\begin{defn}


Пусть $\vp_j\in G^j$ и


$\rho_{j,0}(\vp_j)=(a,b)\in M\times M$.


Символ псевдогруппы $G$ есть


$$


\g^j(\vp_j)=\Ker\big[(\rho_{j,j-1})_*:T_{\vp_j}G^j\to


T_{\vp_{j-1}}G^{j-1}\big]


$$


(здесь можно заменить $G^i$ на $ G^i_{a,b}$). Символ можно


рассматривать как подпространство


$\g^j(\vp_j)\subset S^j(T^*_aM)\ot T_bM\overset{\vp_1^{-1}}{\simeq}{}


S^j(T^*_aM)\ot T_aM$


и отождествлять с алгеброй Ли $\g^j_a$ группы Ли $\frak{G}^j_a$.


\end{defn}





$l$-псевдогруппа $G$ называется {\em формально интегрируемой},


если она $l$-интегрируема, для всех $j>l$ существуют продолжения


$G^j=(G^l)^{(j-l)}$, которые являются псевдогруппами, и проекции


$\rho_{j,j-1}:G^j\to G^{j-1}$ суть векторные расслоения. Также как и в


теории дифференциальных уравнений (\cite{Go}, \cite{Gu}, \cite{S}), критерий


формальной интегрируемости может быть сформулирован в терминах


$\d$-комплекса Спенсера:


$$


%%\begin{equation}\label{hhh1}


0 \to\g^l_a@> \d>>\g^{l-1}_a\ot T^*_aM@>\d>>


\dotsb@>\d>>\g^{l-j}_a\ot\La^jT^*_aM @>\d>> \dotsb


$$


Биградуированные группы когомологий этого комплекса обозначаются через


$H^{l-j,j}(G)$ или через $H^{l-j,j}(\g)$. Препятствиями к формальной


интегрируемости $l$-псевдогруппы $G$, рассматриваемой как


дифференциальное уравнение, являются некоторые элементы $W_j(G)\in


H^{j-1,2}(G)$ (тензоры Вейля или кривизны), которые определяются


геометрией пространств струй (см. \cite{Ly}).








\begin{thm}


Пусть $G$ --- $l$-псевдогруппа. Предположим, что $G$ \ $l$-интегрируема,


символы $\g^j$ над $G^l$ образуют векторное расслоение и все тензоры


Вейля тождественно равны нулю для всех $j\ge l$. Тогда псевдогруппа $G$


формально интегрируема.


\end{thm}





\begin{pf}


Из условий теоремы следует, что $G$ формально интегрируема как


дифференциальное


уравнение (см. \cite{Ly}). Требуется проверить, что полученная система


$\{G^j\}_{j=0}^\infty$ является псевдогруппой, т.\,е. проверить все


условия определения 1. Рассмотрим только первое условие, остальные


проверяются аналогично.





Пусть $G^{j+1}=(G^j)^{(1)}$, $\vp_{j+1}\in G^{j+1}_{a,b}$,


$\psi_{j+1}\in G^{j+1}_{b,c}$ и $\chi_{j+1}=\psi_{j+1}\circ


\vp_{j+1}$. Надо показать, что $\chi_{j+1}\in G^{j+1}$, что


равносильно включению $L(\chi_{j+1})\subset T_{\chi_j}G^j$.


Для того чтобы доказать это включение, рассмотрим


оператор умножения $m_j:G^j\times G^j\to G^j$. Тогда


$T_{\psi_j}G^j\oplus T_{\vp_j}G^j@>{dm_j}>>


T_{\chi_j}G^j$.


Заметим, что прямые слагаемые $T_{\psi_j}G^j$ и $T_{\vp_j}G^j$


содержат подпространства $L(\psi_{j+1})$ и


$L(\vp_{j+1})$ соответственно. Но


$L(\psi_{j+1})\oplus L(\vp_{j+1})@>{dm_j}>>


L(\psi_{j+1}\vp_{j+1})$


для любых $\vp_{j+1},\psi_{j+1}\in D^{j+1}(M)$ таких, что


композиция определена. Отсюда следует требуемое утверждение.


\end{pf}








$l$-псевдогруппа $G$ называется $q${\em -ацикличной}, если


$H^{i,j}(G)=0$ для $i\ge l$,


$0\le j\le q$. \ $\infty$-ацикличная псевдогруппа 


называется {\em инволютивной}.


Для таких псевдогрупп $G$ препятствие к формальной интегрируемости


сводится к единственному препятствию $W_l(G)$.





Если псевдогруппа $G$ формально интегрируема, получаем ее бесконечное


продолжение $G^\infty=\lim_{\op{\,proj}}(G^l,\rho_{l,l-1})$,


которое называется инфинитезимальной или


{\em формальной псевдогруппой}.


Если имеет место локальная интегрируемость, гладкая или


аналитическая (см. Введение), то назовем такую


псевдогруппу {\em интегрируемой}.





\begin{exam}


Группа комплексных дробно-линейных преобразований одномерного


комплексного проективного пространства $S^2=\CC P^1$


(или вещественных дробно-линейных преобразований одномерного


вещественного проективного пространства $S^1=\R P^1$) есть


интегрируемая псевдогруппа порядка 3 и конечного типа. Действительно,


ее алгебра Ли представляет из себя алгебру квадратично-полиномиальных


векторных полей на прямой: $\g=\{\x=(c_0+c_1z+c_2z^2)\p_z\}$.


\end{exam}





\begin{exam}


Псевдогруппа преобразований Ли на пространстве струй


$M=J^k\pi$ для некоторого расслоения $\pi:E_\pi\to M$ (\cite{KLV})


имеет порядок 1 и бесконечный тип. Этот пример является основным для


данной работы и мы обсудим его подробнее в \S\,\ref{S4}.


\end{exam}





\begin{exam}


Пусть $\E$ --- геометрическая структура (\cite{Gu}, \cite{Ly}), и $G$ ---


псевдогруппа Ли ее автоморфизмов. Если структура $\E$ интегрируема


(плоская),


то псевдогруппа тоже интегрируема. Она может иметь конечный или


бесконечный тип в зависимости от $\E$ (\cite{Ko}). Ее порядок


совпадает с порядком структуры $\E$.


Если геометрическая структура не интегрируема, то порядок псевдогруппы


$G$ может возрастать, и она, как правило, не интегрируема (формально


или локально).


\end{exam}








\section{Эквивалентность подмногообразий относительно псевдогрупп}





Возникает естественный вопрос: в каком случае два замкнутых (локальных)


подмногообразия $N_1,N_2\subset M$ эквивалентны относительно


транзитивного действия псевдогруппы $G$? В этом параграфе изучим


инфинитезимальную версию этой проблемы для $l$-струй и $l$-псевдогрупп.





\begin{defn}


Скажем, что $l$-струи двух подмногообразий $N_1$ и $N_2$ в точках


$a,b\in M$ \ $G$-эквивалентны, если $\vp_l[N_1]^l_a=[N_2]^l_b$ для некоторого


$\vp_l\in G^l_{a,b}$.


\end{defn}





Для транзитивных псевдогрупп проблема эквивалентности редуцируется к


случаю, когда $a=b$. В дальнейшем предполагаем, что это равенство


выполнено.


\medskip





2.1. {\it Действие на струях подмногообразий}.


Псевдогруппа $D^l(M)$ и, следовательно, псевдогруппа $G^l$


действуют на пространстве $J^l_r(M)$ по формуле


$\vp_{(l)}:[N]_a^l\mapsto[\vp(N)]_{\vp(a)}^l$. Эти действия


удовлетворяют соотношению


$\rho_{l,s}\circ\vp_{(l)}=\vp_{(s)}\circ\rho_{l,s}$.


Следовательно, группа $\frak{G}^l_a$ действует на $\Fg(a_{l-1})$.


Это действие аффинно для $l>1$ и  порождено линейными


коллинеациями в грассманианах для $l=1$.


Индуцированное (линейное) действие алгебры Ли $\g^l_a\ni\te$ можно


описать следующим образом: $f\mapsto\ll(\te)+f$, $f\in


T_{a_k}\Fg(a_{l-1})$. Здесь


$\ll:S^lT_a^*M\ot T_aM\to T_{a_k}\Fg(a_{l-1})$ отображает элемент


$\te$ в его образ $\ov\te\in S^l\tf^*_a\ot\vg_a$ при ограничении и


последующей факторизации.





Таким образом, стабилизатор элемента $a_l\in\Fg(a_{l-1})$


в случае групп Ли равен


${\frak H}^l_a=\frak{G}^l_a\cap\St$, а в случае алгебр Ли равен


${\frak h}^l_a=\frak{g}^l_a\cap\St$, где


$$


\St=(\Ann\tf_a)\circ_{\op{\,sym}}S^{l-1}T_a^*M\ot


T_aM+S^lT^*_aM\ot\tf_a.


$$


В частности, т.\,к. символ псевдогруппы $D^l(M)$ действует


транзитивно, получаем


\begin{equation}\label{zxc}


S^lT^*_aM\ot T_aM/\St\simeq S^l\tf^*_a\ot\vg_a.


\end{equation}





\begin{note}\label{rk1}


Так как группа $\frak{G}^l_a$ коммутативна при $l>1$, то нет различия между


действиями групп Ли и алгебр Ли. В случае 1-струй действия различны.


\end{note}





Рассмотрим далее некоторые выделенные подмногообразия коразмерности  $r$


в $M$. Выделение подмножества таких подмногообразий из множества всех


подмногообразий зависит от выбора некоторой общей внутренней


структуры относительно действия псевдогруппы и будет описано ниже.


Оно задается некоторыми дифференциальными соотношениями, поэтому


рассмотрим множество струй подмногообразий, выделенных некоторым


дифференциальным уравнением $\frak{N}$.





Итак, пусть $\frak{N}\subset J^l_r(M)$~--- $G$-инвариантное


дифференциальное уравнение. Как обычно, его символ


$h_a^l\subset S^l\tf^*_a\ot\vg_a$ есть $\rho_{l,l-1}$-вертикальное


подпространство


в $T_{a_l}\frak{N}$. Так как псевдогруппа $G$ действует на


$\frak{N}$, получаем следующую точную последовательность:


\begin{equation}\label{130300}


0\to{\frak h}^l_a\hookrightarrow


\g^l_a@>{\ll}>> h_a^l@>{\varpi}>> {\frak O}^l_a\to 0.


\end{equation}





\begin{defn}


Фактор ${\frak O}^l_a=h_a^l/\ll(\g^l_a)$ называется пространством


$l$-{\em ковариантов\/} действия псевдогруппы $G$. Двойственное


пространство $({\frak O}^l_a)^*$ называется пространством


$l$-{\em вариантов}.


\end{defn}





Изучение формальной эквивалентности подмногообразий относительно


$G$-действия будем проводить индуктивно, опираясь на следующее


очевидное утверждение.





\begin{prop}


Пусть $[N_1]^{l-1}_a=[N_2]^{l-1}_a\in\rho_{l,l-1}(\frak{N})$ и


$l>1$. \ $l$-струи подмногообразий $N_1$ и $N_2$ из $\frak{N}$ в


точке $a\in M$ \ $G$-эквивалентны тогда и только тогда, когда


они принадлежат одной и той же $\g^l_a$-орбите в $h_a^l$, которая


является аффинным подпространством коразмерности равной


$\dim{\frak O}_a^l=\dim h_a^l- \dim\big(\g^l_a/{\frak


h}^l_a\big)$, иначе говоря, тогда и только тогда, когда


ростки подмногообразий имеют одинаковые $l$-варианты\/${:}$


$\varpi([N_2]_a^l-[N_1]_a^l)=0$. 


\end{prop}





Требование $l>1$ объясняется замечанием \ref{rk1}. При $l=1$


действия групп Ли и алгебр Ли различны: орбиты групп Ли


лежат в $\Grass_r(T_aM)$, в то время как орбиты алгебр Ли


являются аффинными подпространствами в касательном пространстве к


грассманиану в $a_1$. Поэтому 1-струи требуют отдельного рассмотрения.


\medskip





2.2. {\it Дифференциальные инварианты и эквивалентность}.


Пусть $\cal{I}$ есть алгебра дифференциальных инвариантов


псевдогруппы Ли $G$, т.\,е. функций постоянных на орбитах


$G$-действия на $\frak{N}$. Обозначим через $\cal{I}_k$


подалгебру инвариантов порядка $k\le l$ (поднятых с пространства


$k$-струй).


Зафиксируем точку $a_l\in J^l_r(M)$ и определим возрастающую


фильтрацию пространства $T_{a_l}^*J^l_r(M)$ подпространствами


$$


\T_k(a_l)=\{d_{a_l}f\mid f\in\cal{I}_k\}\subset


T_{a_l}^*J^l_r(M),\quad k=0,\dotsc,l.


$$





Заметим, что $\T_l$~--- уравнение первого порядка, определяющее


$G^l$-дифференциальные инварианты на $J^l_r(M)$. Около сингулярных


орбит дифференциальные инварианты устроены плохо, поэтому в них


определяем фильтрацию следующим образом (в регулярных точках


определения совпадают):


$\T_k(a_l)=\pi_{l,k}^*\Ann T_{a_k}(G^k\cdot a_k)$.








\begin{prop}


При $0<k\le l$ выполняется $\frak{O}^k_a=(\T_k/\T_{k-1})^*$.


\end{prop}





\begin{pf}


Действительно, $\frak{O}^k_a=T_{a_k}(\pi_{k,k-1})^{-1}(G^{k-1}\cdot


a_{k-1})/T_{a_k}(G^k\cdot a_k)$, откуда следует утверждение.


\end{pf}





Таким образом, можно получить решение проблемы формальной


эквивалентности с помощью следующей индуктивной процедуры.


Начнем с псевдогруппы $G$ и всех подмногообразий, т.\,е. с


$\frak{N}=J^l_r(M)$. Пусть $\frak{O}^l$~--- первое нетривиальное пространство


$l$-ковариантов. Для того чтобы задать в нем точку, зафиксируем


$l$-варианты в $(\frak{O}^l_a)^*=\T^l$, что эквивалентно фиксации


значений дифференциальных инвариантов порядка $l$. Другими словами,


фиксируется тип внутренней геометрии подмногообразий. Отсюда


получается меньшее уравнение ${\frak N}\subset J^l_r(M)$ на


подмногообразия, и процедура повторяется.


В регулярных точках процедура завершится за конечное число шагов в


силу теоремы о продолжении Картана--Кураниши.





\begin{note}


Последнее утверждение эквивалентно тому, что пространство


дифференциальных инвариантов конечно порождено относительно линейных


комбинаций и производных Трессе.


\end{note}








2.3. {\it Транзитивность и трансверсальность}.


Теперь требуется найти критерий, позволяющий проверять отсутствие


$l$-вариантов, или эффективный метод вычисления дифференциальных


инвариантов.





\begin{defn}\label{dfn5}


Говорят, что псевдогруппа $G$ действует {\em $l$-транзитивно}


около $a_l\in\frak{N}$, если для любой другой струи $b_l\in\frak{N}$,


близкой к $a_l$, существует струя $\vp_l\in G^l_{a,b}$ такая, что


$\vp_l(a_l)=b_l$. Другими словами, орбита $G^l\cdot a_l$ открыта.


\end{defn}





\begin{defn}


Скажем, что действие псевдогруппы $G$ \ {\em $l$-трансверсально}


около $a_l$, если требование определения \ref{dfn5} выполняется


при условии $a_{l-1}=b_{l-1}$. Другими словами,


$\frak{G}^l_a$ действует транзитивно на $\frak{F}(a_{l-1})$.


\end{defn}





Для того чтобы объяснить эту терминологию, рассмотрим отображение


$\ll:\te\mapsto\ov\te$ из п.\,2.1.


Пространство $\ll^{-1}(h^l_a)\subset S^lT^*_aM\ot


T_aM$ содержит два подпространства $\St$ и $\frak{g}^l_a$.





Пусть $l>1$. Следующее утверждение следует из (\ref{zxc}),


(\ref{130300}) и определений.





\begin{prop}\label{thm4}


$l$-трансверсальность $G$ на $\frak{N}$ эквивалентна любому из


следующих условий\/${:}$


\begin{itemize}


\item[---] $\St$ трансверсально $\frak{g}_a^l$ в $\ll^{-1}(h_a^l){:}$


$\St+\frak{g}_a^l=\ll^{-1}(h_a^l)$,


\item[---] нет $l$-ковариантов, т.\,е. $\frak{O}_a^l=0$.


\end{itemize}


\end{prop}





Ясно, что $l$-трансверсальность есть индуктивный шаг для получения


$l$-транзитивности (1-струи $a_1$, $b_1$ надо рассматривать отдельно


согласно замечанию \ref{rk1}).





\begin{thm}


Пусть $G^1\cdot a_1$ открыто и $G$ действует $j$-трансверсально на


$a_j$ для $1<j\le l$. Тогда $G$ действует $l$-транзитивно около $a_l$.


\end{thm}








\section{Основные примеры}





В этом разделе проверяются условия трансверсальности для псевдогрупп


автоморфизмов ряда основных геометрических структур.


Неприводимые псевдогруппы Ли были классифицированы Э.\,Картаном


(\cite{C}). Для псевдогруппы


$G=\Diff_{\loc}(M)$ всех диффеоморфизмов и псевдогруппы


$G=\Diff_{\loc}(M,\Omega)$ диффеоморфизмов, сохраняющих


объем, все подмногообразия коразмерности


$r$ локально $G$-эквивалентны. 





Рассмотрим другие случаи.


Заметим, что для псевдогруппы преобразований некоторой неинтегрируемой


структуры $l$-псевдогруппа $G^l$ состоит из струй диффеоморфизмов,


сохраняющих эту структуру до порядка $l$.


Поэтому, несмотря на обозначение, при $j<\min(k,l)$ подпсевдогруппа


$G^j$ может быть разной в зависимости от вложения в $G^k$ или в $G^l$.


\medskip





3.1. {\it Псевдогруппа преобразований комплексной структуры}.


Пусть $G$ --- псевдогруппа локальных голоморфных преобразований


комплексного многообразия $(M,J)$ комплексной размерности~$n$. Тогда


символ --- $\g^l_a=S^lT^*_aM\ot_\CC T_aM$.





\begin{prop}\label{CC}


Псевдогруппа $G$ \ $l$-трансверсальна около $a_l=[N]_a^l$ тогда и


только тогда, когда


\begin{itemize}


\item[1.]


$\tf_a\cap J\tf_a=\{0\}$ или $\tf_a+J\tf_a=T_aM$ при $l=1;$


\item[2.]


$\tf_a\cap J\tf_a=\{0\}$ при $l>1$.


\end{itemize}


\end{prop}





\begin{pf}


{\bf Необходимость.} Условие 1 означает, что не могут одновременно


существовать нетривиальные комплексные подпространства $L'\subset\tf_a$,


$L''\cap\tf_a=\{0\}$. Если это неверно, то образ


$\ll:\g_a^1\to h_a^1=\tf_a^*\ot\vg_a$


состоит из отображений $A:\tf_a\to\vg_a$


таких, что индуцированное отображение $\ti A:L'\to L''$ линейно над


полем комплексных чисел:


$\ti A(J\x)=J\ti A(\x)$. Следовательно, существуют тензоры в $h_a^1$,


не лежащие в образе $\ll$. Аналогично, при $l>1$, если существует


комплексное подпространство $L'\subset\tf_a$, то ограничение каждого


$A\in\ll(\g_a^l)$ на него удовлетворяет условию


$\ti A(\x_1,\x_2,\x_3,\dotsc)+ \ti


A(J\x_1,J\x_2,\x_3,\dotsc)=0$ и, следовательно, не произвольно.


\smallskip





{\bf Достаточность.} Если $J\tf_a\cap\tf_a=\{0\}$, то каждый тензор


$R\in S^l\tf^*_a\ot\vg_a$ может быть продолжен до тензора


$\wh R\in S^lT^*_aM\ot_\CC T_aM$. Действительно, разложим


$T_aM=\tf_a\oplus J\tf_a\oplus W$, где $JW=W$, и положим


$\wh R(J\x,\cdot)=J\wh R(\x,\cdot)$, $\x\in\tf_a$, и $\wh


R(w,\cdot)=0$, $w\in W$.


Аналогично, можно показать, что если $\tf_a+J\tf_a=T_aM$, то имеет


место 1-трансверсальность.


\end{pf}








В частности, струя подмногообразия общего положения размерности


$\dim_{\R} N\le n$ \ $l$-транс\-вер\-саль\-на для любого $l\ge0$, поэтому


комплексная псевдогруппа $G$ действует на таких локальных


(аналитических) подмногообразиях транзитивно. С другой стороны,


подмногообразие размерности $\dim_{\R} N>n$ не может быть


трансверсальным.


Действительно,  $\Pi_a=\tf_a\cap J\tf_a\ne\{0\}$,


следовательно, $N$ обладает внутренней геометрией.





Изучение многообразий $N$, наделенных распределением с комплексной


структурой, является предметом $CR$-геометрии. Ее инвариантами являются


кривизны Картана--Чженя--Мозера \cite{CM}.


Если зафиксировать эти кривизны, получим меньший класс $\frak{N}$


подмногообразий, на котором действие трансверсально, и, следовательно,


транзитивно.





Другой важный класс $\frak{N}$ состоит из всех комплексных


подмногообразий $N\subset M$ \ $\CC$-ко\-раз\-мер\-но\-сти $r$. Этот класс


$l$-трансверсален для каждого $l$, и поэтому транзитивен.


\medskip





3.2. {\it Псевдогруппа преобразований почти комплексной структуры}.


Pассмотрим случай неинтегрируемой почти комплексной структуры


$J$, $J^2=-\bold 1$. Псевдогруппа $G^l$ состоит из всех


$J$-голоморфных $l$-струй: $J\circ d\vp_l=d\vp_l\circ J$.


Если $l=1$, то $G^1_a=\GL_\CC(T_aM)$, как и в комплексном


случае.





Пусть $N_J\in\Hom_{\ov\CC}(\La^2TM,TM)$~--- тензор Нейенхейса


структуры $J$ (напомним, что он является препятствием к


интегрируемости). При $l=2$ имеем


$$


G^1_a=\rho_{2,1}(G^2_a)=\{\Phi\in T^*_aM\ot T_aM \mid


J\c\Phi=\Phi\c J,\ N_J\c(\Phi\we\Phi)=\Phi\c N_J \}.


$$


Симметрическая связность без кручения $\nabla$ задает разложение


2-струи $\vp_2\in G^2$ с компонентами $(a,\Phi,\Phi^{(2)})$.


Для $\vp_1=(a,\Phi)$ последняя компонента $\Phi^{(2)}\in\Fg(\vp_1)$


задается следующим образом:


$$


\{\Phi^{(2)}\in S^2T^*_aM\ot T_aM,\


J\P^{(2)}(\x,\e)-\P^{(2)}(J\x,\e)=\P\c\nabla_\e


(J)(\x)-\nabla_{\P\e}(J)(\P\x)\},


$$


Таким образом, $\g_a^2=S^2T^*_aM\ot_\CC T_aM$ как и в комплексном


случае, но для меньшего множества точек $\vp_1$. \ 2-псевдогруппа


$G^2$, вообще говоря, не является 2-интегрируемой. Доказательство этих


фактов и описание проекции $\rho_{l,l-1}:G^l_a\to G^{l-1}_a$


приведено в \cite{Kr1}.





Можно показать (см. \cite{Kr2}), что для структуры $J$ общего


положения множество $G^2$ состоит только из единицы для $n>3$, множество $G^3$


состоит только из единицы для $n>2$ и множество $G^4$ --- из единицы для $n=2$


(здесь не рассматривается случай $n=1$, при котором $J$ всегда


интегрируема). Анализ псевдоголоморфных инвариантов струй


подмногообразий на основе классификации тензоров Нейенхейса (см.


\cite{Kr2}) приводит к следующему результату.





\begin{prop}


Пусть $(M,J)$~--- почти комплексное многообразие с неинтегрируемой


структурой $J$ общего положения и $n>1$. Для $l=1$ трансверсальность


описывается условием $1$ предложения $\ref{CC}$. Для $l=2$ никакая


$2$-струя не является трансверсальной, кроме случая $n=2$ и $\dim_{\R}


N=1$. Трансверсальность также отсутствует при $l \ge 3$.


\end{prop}








3.3. {\it Псевдогруппа преобразований римановой метрики}.


Вначале рассмотрим псевдогруппу изометрий евклидова пространства


$\R^n$. Ее глобализация есть группа $G=O(n)\leftthreetimes\R^n$.


Эта псевдогруппа имеет конечный тип и $\g^l_a=0$ при $l\ge2$


(\cite{Ko}). Поэтому трансверсальности нет для $l\ge2$, но действие


является 1-трансверсальным около каждой 1-струи $\vp_1$.





Теперь рассмотрим риманово многообразие $(M^n,q)$ и пусть $G$~---


псевдогруппа изометрий. При $l=1,2$ группа $G^l$ такая же как в


евклидовом случае. Рассмотрим $l=3$. Тогда $G^1_a=\rho_{3,1}(G^3_a)$


состоит из линейных изометрий из $O(T_aM)$, сохраняющих риманову


кривизну.





Как и в почти комплексном случае, для структуры $q$ общего положения


псевдогруппа $G^l$ состоит только из единицы при $l>3$ или


$l=3$, $n>2$.





\begin{prop}


При $l>1$ не существует $1$-струи, около которой действие $G$


трансверсально.


\end{prop}





Действительно, различные внутренние и внешние кривизны являются


$l$-вариантами.


Фиксируя их, получаем трансверсальность и отсюда эквивалентность.


\medskip





3.4. {\it Псевдогруппа симплектических преобразований}.


Рассмотрим симплектическое многообразие $(M,\oo)$ размерности $2n$ и


пусть $G$~--- (псевдо)группа его симплектоморфизмов. Используя


изоморфизм $TM\overset{\oo}{\simeq}{} T^*M$, запишем символы в виде


$\g^l_a=S^{l+1}T^*_aM\subset S^lT^*_aM\ot T^*_aM$, отождествляя их


с однородными порождающими функциями (гамильтонианами) степени $l+1$.





\begin{prop}\label{symplpro}


$G$ действует $l$-трансверсально около $a_l\in J^l_r(M)$ для всех


$l\ge1$ тогда и только тогда, когда ограничение $\oo$ на $\tf_a=a_1$


имеет максимальный ранг.


\end{prop}





\begin{pf}


Необходимость этого условия очевидна. Для того чтобы доказать


утверждение в обратную сторону, рассмотрим вначале случай $\dim\tf_a\in2\Z$.


Запишем разложение в прямую сумму $TM=\tf\oplus\vg$, где слагаемые


симплектически ортогональны. Условие трансверсальности в силу


предложения \ref{thm4} записывается в виде


$$


\vg^*\circ S^{l-1}T^*M\ot T^*M+ S^lT^*M\ot\tf^*+ S^{l+1}T^*M = S^lT^*M\ot


T^*M,


$$


которое очевидно выполняется.





Для $\dim\tf_a\in2\Z+1$ запишем разложение


$TM=\frak{u}\oplus(\frak{l}\oplus\frak{r})\oplus\frak{v}$


в симплектически ортогональную сумму,


где $\tf=\frak{u}\oplus\frak{l}$ и $\frak{u}$,


$\frak{v}$ симплектические, $\dim\frak{l}=\frak{r}=1$.


Теперь трансверсальность


$$


(\frak{r}^*\oplus\frak{v}^*)\circ S^{l-1}T^*M\ot T^*M+


S^lT^*M\ot(\frak{u}^*\oplus\frak{r}^*)+ S^{l+1}T^*M = S^lT^*M\ot T^*M


$$


следует из тождества


$S^i\frak{l}^*\ot\frak{l}^*=S^{i+1}\frak{l}^*$.


\end{pf}








Полученное утверждение эквивалентно частному случаю теоремы


Вайнштейна--Гивенталя (\cite{AG}).


Для того чтобы исследовать более общий случай, надо рассмотреть


различные ранги ограничения $\oo|_N$. Тогда


трансверсальность уже не имеет места и получаем 1-вариант, который


очевидно есть ранг (или размерность


$\Ker(\oo|_N)$). Зафиксировав его, получим трансверсальность для 


соответствующего уравнения $\frak{N}$ на подмногообразие.





Наконец, рассмотрим класс $\frak{N}$ изотропных или коизотропных


подмногообразий.


Аналогичные вычисления показывают, что $G$ действует на них


$l$-трансверсально для любого $l$.


\medskip





3.5. {\it Псевдогруппа контактных преобразований}.


Рассмотрим контактное многообразие $(M,\Pi^{2n})$, $\dim M=2n+1$,


и обозначим через $\nu=TM/\Pi$ его нормальное расслоение. Пусть $G$~---


(псевдо)группа контактных преобразований. Ее алгебра Ли состоит из


контактных векторных полей $X_f$, которые определяются порождающими


сечениями (гамильтонианами) $f\in C^\infty(M)\ot\nu$.





Выбор ненулевого сечения расслоения $\nu$ эквивалентен выбору


контактной формы $\a\in C^\infty(\Ann(\Pi)\setminus 0)$,


$\a\we d\a^n\ne0$. Тогда гамильтониан можно считать функцией,


$f\in C^\infty(M)$, и контактное поле единственным образом задается


равенствами


$$


\a(X_f)=f,\ \ d\a(\cdot,X_f)=df|_{\Pi}.


$$


В координатах Дарбу $(q,u,p)$, $\a=du-p_idq^i$, имеем


$$


X_f=\D_{q^i}(f)\p_{p_i}-\p_{p_i}(f)\D_{q^i}+f\p_u, \quad


\text{где }\ \D_{q^i}=\p_{q^i}+p_i\p_u.


$$





Заметим, что задание $\a$ эквивалентно расщеплению


$T_aM=\Pi_a\oplus\nu_a$, где первое слагаемое~--- симплектическое


подпространство, а второе слагаемое --- одномерное евклидово


пространство.


Для того чтобы описать символ $\g^l_a$, отождествим $TM\simeq T^*M$


по прямым слагаемым, используя симплектическую структуру на $\Pi$ и


евклидову структуру на $\nu$. Тогда получим


$$


\g^l\simeq S^l\nu^*\oplus\sum_{i>0} S^i\Pi^*\ot


S^{l+1-i}\nu^*\simeq S^{l+1}T^*M.


$$


Действительно, контактные векторные поля $X_f$ порядка $l$


определяются гамильтонианами $f$ порядка $(l+1)$ по всем переменным,


кроме чисто $l$-й степени переменной $u$.





\begin{prop}


Действие $G$ является $l$-трансверсальным около $\vp_l=[N]_a^l$ для


всех $l\ge1$ тогда и только тогда, когда $\tf_a=T_aN$ трансверсально


контактной плоскости $\Pi_a$ и структура на $\Pi_a^N=\Pi_a\cap T_aN$,


индуцированная из канонической конформно-симплектической структуры на


$\Pi$, максимально невырождена. Это значит, что


\begin{itemize}


\item[---] если $\dim\Pi^N_a=2r$, то


$(d\a)^r|_{\Pi^N_a}\ne0$,


\item[---] если $\dim\Pi^N_a=2r+1$, то


$\rank(d\a|_{\Pi^N_a})=2r$.


\end{itemize}





Эти условия эквивалентны утверждению, что через любую точку близкую к


$a\in N$ проходит изотропное подмногообразие размерности не


превосходящей $r$.


\end{prop}





Так как символы те же, что и в симплектическом случае, доказательство


можно провести аналогично доказательству предложения \ref{symplpro}.


Другой способ найти символ $\g^l$ и доказать утверждение состоит в


использовании естественной фильтрации пространства $S^lT^*M\ot TM$


степенями $\Ann(\Pi)$ в первом сомножителе и степенями $\Pi$ во


втором.


%%% \end{pf}





Как в симплектическом случае, заметим, что


$\rank((d\a)^r|_{\Pi^N_a})$


есть 1-вариант, фиксируя который получаем


трансверсальность. Имеет место контактный аналог теоремы


Вайнштейна--Гивенталя, который стандартно доказывается с помощью


изотопии. Для простоты привeдeм здесь только локальную версию (хотя


глобальная тоже верна, если наложить дополнительное условие на


нормальное расслоение).





\begin{thm}


Пусть $N_1$, $N_2$~--- локальные подмногообразия коразмерности $r$ в


контактном


многообразии $(M,\Pi)$. Если ограничения контактной структуры $\Pi$ на


эти подмногообразия изоморфны, то у них есть контактно изоморфные


окрестности. 


\end{thm}





Наконец, аналогично п.\,3.4, как частный случай, получаем, что в


ограничении на класс $\frak{N}$ изотропных подмногообразий фиксированной


размерности, имеет место трансверсальность $G$-действия.








\section{Псевдогруппа преобразований Ли}\label{S4}





Пусть $\pi$~--- векторное расслоение и $G$~--- псевдогруппа


преобразований Ли пространства $M=J^k\pi$. Она состоит из локальных


диффеоморфизмов расслоения струй, сохраняющих распределение Картана


$\C_k$ (\cite{KLV}).


Эквивалентно, преобразование Ли сохраняет класс изотропных


подмногообразий метасимплектической структуры. Теорема Ли--Бэклунда


(\cite{KLV})


утверждает, что преобразование Ли есть поднятие диффеоморфизма пространства


$J^0\pi$, если $r:=\rank\pi>1$ или поднятие контактного


преобразования пространства $J^1\pi$, если $r=1$. Поэтому


преобразования Ли являются поднятиями преобразований пространства


$J^\varepsilon\pi$, где $\varepsilon=\max(0,2-r)$.





Псевдогруппа $G$ интегрируема. Ниже мы найдем ее символы.





Векторное поле называется инфинитезимальным преобразованием Ли, если


его поток состоит из преобразований Ли. Для того чтобы описать эти поля


Ли, выберем координатную окрестность $(x,u)$ на расслоении $\pi$


(бескоординатное описание см. в \cite{KLV}). Она индуцирует


координаты $(x^i,p^j_\z)_{0\le|\z|\le k}$ на


$J^k\pi$, где $p^j_\z([u]_x^k)=\frac{\p^{|\z|}u^j}{\p x^\z}$.





Напомним, что оператор полной производной


$D_i=\p_{x^i}+\sum\limits_{j;\z}p^j_{\z+1_i}\p_{p_\z^j}$ отображает


$C^\infty(J^k\pi)$ в $C^\infty(J^{k+1}\pi)$. Полная производная,


соответствующая мультииндексу $\z=(i_1,\dots,i_s)$, обозначается через


$D_\z=D_{i_1}\circ\dotsb\circ D_{i_s}$.





Обозначим через $\t_x$ пространство, касательное к базе расслоения


$\pi$ в точке $x=\pi_k(x_k)$ и через $\nu_{x_0}$ пространство,


касательное к слою в $x_0=\pi_{k,0}(x_k)$. Также пусть  $F(x_k)$~---


$\pi_{k+1,k}$-слой и $\ups_{x_1}=T_{x_1}\big(F(x_0)\big)$.


\medskip





4.1. {\it Точечные преобразования высшего порядка}.


Для $r>1$ обозначим проекцию поля Ли на


$J^\varepsilon\pi=J^0\pi$ через


$X=\sum\limits_ia^i(x,u)\p_{x^i}+\sum\limits_jb^j(x,u)\p_{u^j}$.


Тогда продолжение на $J^k\pi$ имеет вид


\begin{equation}\label{X^k}


X^{(k)}=\sum_ia^i(x,u)D_i^{(k+1)}+\sum_{j;|\z|\le k}


D_\z(\vp^j)\p_{p^j_\z},


\end{equation}


где $\vp^j=b^j-\sum\limits^n_{i=1}a^ip^j_i$~--- компоненты так называемой


{\em производящей функции} $\vp=(\vp^1,\dotsc,\vp^r)$,


и $D_i^{(k+1)}=\p_{x^i}+\sum\limits_{j;|\z|\le


k}p^j_{\z+1_i}\p_{p_\z^j}$ есть оператор полной производной,


ограниченной на $J^k\pi$. Хотя кажется, что коэффициенты (\ref{X^k})


зависят от $(k+1)$-струй, поле Ли на самом деле определено на $J^k\pi$.





Формула (\ref{X^k}) следует из того, что поле Ли сохраняет


кораспределение


$$


\Ann(\C_k)=\Big\langle\oo_\z^j=dp_\z^j-\sum_ip_{\z+1_i}^jdx^i


\,\big|\, 1\le j\le r,\ |\z|<k\Big\rangle


$$


и из того, что $d=\sum\limits_i dx^i\ot D_i^{(k+1)}+ \sum\limits_{j;|\z|\le


k}\oo_\z^j\ot\p_{p_\z^j}$ на $J^k\pi$.





\begin{prop}\label{proppre}


$l$-символ $\g^l(x_k)$ псевдогруппы $G$ в точке


$x_k\in J^k\pi$ допускает расщепление $\g^l=\g^l_H\oplus\g^l_V$,


зависящее от точки $x_{k+1}\in F(x_k)$. Горизонтальная компонента


изоморфна


$$


\g^l_H(x_k)\simeq


\bigg[S^l\nu_{x_0}^*\oplus\sum_{0<i<k}(S^i\tau_x^*\ot


S^{l-1}\nu_{x_0}^*)\oplus\sum_{i\ge k}(S^i\tau_x^*\ot


S^{k+l-i-1}\nu_{x_0}^*)\bigg]\ot\tau_x,


$$


в то время как вертикальная $($эволюционная$)$ компонента представляется


в виде


$$


\g^l_V(x_k)\simeq \bigg[\sum_{0\le i<k}(S^i\tau_x^*\ot


S^l\nu_{x_0}^*)\oplus\sum_{i\ge k}(S^i\tau_x^*\ot


S^{k+l-i}\nu_{x_0}^*)\bigg]\ot\nu_{x_0}.


$$


\end{prop}





\begin{pf}


Пространство $T_{x_k}J^k\pi$ разлагается в прямую сумму


горизонтальной\linebreak $L(x_{k+1})\subset\C_k(x_k)$ и вертикальной


$T_{x_k}^v=\Ker(\pi_k)_*$ компонент. Отсюда получаем, что


$$


\g^l(x_k)\subset S^lT^*_{x_k}J^k\pi\ot T_{x_k}J^k\pi =


\bigl[S^lT^*_{x_k}J^k\pi\ot L(x_{k+1})\bigr]\oplus


\bigl[S^lT^*_{x_k}J^k\pi\ot T_{x_k}^vJ^k\pi\bigr],


$$


что влечет требуемое расщепление. В формуле (\ref{X^k})


горизонтальная (вертикальная) компонента соответствует первому (второму)


слагаемому.





Обозначим через $\mu_a$ идеал в $C^\infty(M)$, состоящий из функций,


обращающихся в нуль в $a\in M$, и через $\mu^l_a$ его степень. Пусть


$\mu_a^l(\frak{Lie})$ --- пространство полей Ли, имеющих в $a$


нуль порядка $l$. Тогда


$\g^l(x_k)=\mu_{x_k}^l(\frak{Lie})/\mu_{x_k}^{l+1}(\frak{Lie})$.





Как в контактном, так и в симплектическом случаях мы представляем символ с


помощью струй производящих функций. Он вкладывается в пространство


$S^lT^*_{x_k}J^k\pi\ot T_{x_k}J^k\pi$ в силу (\ref{X^k}).





Выберем координатную систему так, чтобы точка $x_k$ была


ее центром. Если $x_k=[s]^k_a$ для некоторого сечения $s$, то этого


можно добиться, выбирая $s$ за нулевое сечение: $s=\{u^j=0\}$. Тогда


условие $X^{(k)}\in\mu^l_{x_k}$ следующим образом выражается в терминах


компонент производящей функции


$$


a^i\in\mu^l_{x_0},\ \ \p_{x^\z}(a^i)\in\mu^{l-1}_{x_0},\quad


b^j\in\mu^l_{x_0},\ \ \p_{x^\z}(b^j)\in\mu^l_{x_0},\quad


0\le|\z|\le k.


$$


Это доказывает требуемое утверждение. Заметим, что разложение


$T_{x_0}^*J^0\pi=\tau_x^*\oplus\nu_{x_0}^*$ индуцировано точкой


$x_1$, поэтому доказанное представление является каноническим.


\end{pf}








4.2. {\it Контактные преобразования высшего порядка}.


Для $r=1$ преобразование Ли определяется контактным преобразованием


$X^{(1)}=X_\vp$ на $J^1\pi$, соответствующим производящей функции


$\vp=\vp(x^i,u,p_i)$:


$$


X^{(1)}=\sum_i\big[D_i^{(1)}(\vp)\p_{p_i}-\p_{p_i}(\vp)D_i^{(1)}\big]


+\vp\p_u.


$$


Продолжение этого поля на $J^k\pi$ задается формулой,


аналогичной (\ref{X^k}),


\begin{equation}\label{X^kk}


X^{(k)}=-\sum_i \p_{p_i}(\vp)D_i^{(k+1)}+\sum_{|\z|\le


k}D_\z^{(k)}(\vp)\p_{p_\z}.


\end{equation}


Вновь вычисление показывает, что это есть векторное поле на $J^k\pi$,


совпадающее с $X_\vp$ при $k=1$.





Далее используем разложение


$T_{x_1}J^1\pi=\tau_x\oplus\nu_{x_0}\oplus\ups_{x_1}$, которое уже не


является каноническим. Хотя точка $x_2$ определяет расщепление


$T_{x_1}J^1\pi=L(x_1)\oplus T^v_{x_1}$, второе слагаемое далее


раскладывается с помощью связности в расслоении $\pi_{1,0}$.





\begin{prop}\label{proppro}


$l$-символ псевдогруппы $G$ в точке $x_k\in J^k\pi$ есть


\begin{multline*}


\g^l(x_k)\simeq \sum_{0\le j\le l}(S^j\nu_{x_0}^*\ot


S^{l+1-j}\ups_{x_1}^*)\oplus S^l\nu_{x_0}^*\oplus\sum_{1\le


i<k;j}(S^i\tau_x^*\ot S^j\nu_{x_0}^*\ot S^{l-j}\ups_{x_1}^*)\oplus\\


\oplus \sum_{k\le i;j}(S^i\tau_x^*\ot S^j\nu_{x_0}^*\ot


S^{k+l-i-j}\ups_{x_1}^*).


\end{multline*}


\end{prop}





\begin{pf}


Как и при доказательстве предложения \ref{proppre}, в силу (\ref{X^kk})


относительно системы координат $(x^i,u)$, в которой


$p_\z(x_k)=0$ при $|\z|\le k$, условие


$X^{(k)}\in\mu^l_{x_k}(\frak{Lie})$ эквивалентно условиям


$$


\vp\in\mu^l_{x_1},\ \ \p_{p_i}(\vp)\in\mu^l_{x_1},\ \


\p_{x^\z}(\vp)\in\mu^l_{x_1},\quad 0\le|\z|\le k,


$$


откуда и следует утверждение.


\end{pf}





\begin{note}


Другое представление символа преобразований Ли, отличное от


приведенного в предложениях \ref{proppre}--\ref{proppro}, можно


получить с помощью естественной фильтрации пространства


$S^lT^*_{x_k}J^k\pi\ot T_{x_k}J^k\pi$ \ $\pi_{j,j-1}$-проекциями.


\end{note}








4.3. {\it Дифференциальные уравнения}.


Подмногообразие в $J^k\pi$ может рассматриваться как (регулярное)


дифференциальное уравнение, если оно является подрасслоением


относительно всех $\pi_{j,j-1}$-проекций. В локальной ситуации этот


случай является случаем общего положения.





Рассмотрим два таких дифференциальных уравнения: $\E\subset J^k\pi$ и


$\E'\subset J^k\pi'$, а также зададим две точки $x_k$, $x_k'$. Тогда


существует преобразование Ли $\vp:J^\varepsilon\pi\to J^\varepsilon\pi'$,


для которого $\vp^{(k)}(x_k)=x_k'$. Таким образом, мы свели задачу к


случаю, когда оба дифференциальных уравнения $\E$ и


$\E'_\vp=\left(\vp^{-1}\right)^{(k)}(\E')$ определены на одном


пространстве $J^k\pi$. Теперь попытаемся отождествить $\E$ с


$\E'_\vp$ с помощью преобразования Ли. Соответственно получим


точечную или контактную эквивалентность дифференциальных уравнений.





Так получаются внешние симметрии. Есть еще внутренние симметрии,


которые во многих случаях совпадают с внешними (жесткие и, в


частности, нормальные уравнения \cite{KLV}). В этой статье мы не


будем их рассматривать.








\section{Размерностное препятствие к трансверсальности}





В силу предложения~\ref{thm4} условие $l$-трансверсальности влечет


следующее неравенство для символа $h_a^l$ рассматриваемого класса


подмногообразий $\frak{N}$


(дифференциального уравнения):


\begin{equation}


\dim\g^l_a\ge\dim h_a^l. \label{reform-trans}


\end{equation}


Это легко проверяемое условие часто бывает полезно. В следующих


примерах показывается, что его выполнение влечет трансверсальность для


подмногообразия $N$ {\em общего положения}.





В этом разделе рассматриваются произвольные подмногообразия $N\subset M$


коразмерности $r$. В частности, $h_a^l=S^l\tf^*_a\ot\vg_a$.


\medskip





5.1. {\it Вычисления для основных примеров}.


1. Пусть $M$~--- комплексное многообразие размерности


$\dim_\CC M=n$ и $G$~--- псевдогруппа из п.\,3.1. В этом


случае условие (\ref{reform-trans}) записывается в виде


$$


\tbinom{l+n-1}{n-1}\cdot2n\ge\tbinom{l+2n-r-1}l\cdot r.


$$


Это неравенство выполняется при $r\ge n$, а при $r<n$ оно неверно,


если $l>1$. Тем не менее, заметим, что для $l=1$ данное неравенство


имеет место для всех $0<r<n$. В этом случае $(\ref{reform-trans})$ не


достаточно для 1-трансверсальности всех $N$, но достаточно для


подмногообразий $N$ общего положения в точке $a$.





2. В случае почти комплексной или римановой структуры общего


положения имеем $\g^l=0$ для больших $l$, следовательно, трансверсальность


не имеет места.





3. Для псевдогрупп симплектических и контактных преобразований


$\g^l\simeq S^{l+1}T^*M$. Поэтому, если $M=n$, условие


(\ref{reform-trans}) принимает вид


$$


\tbinom{l+n}{n-1}\ge\tbinom{l+n-r-1}{n-r-1}\cdot r,


$$


что всегда верно. Следовательно, никакого условия на размерность


подмногообразия нет и подмногообразие $N\subset M$


общего положения $l$-трансверсально для любого $l$.


\medskip





5.2. {\it Псевдогруппа Ли точечных преобразований}.


Вычислим ограничения на размерность для псевдогруппы преобразований Ли.


Начнем с преобразований, поднятых с $J^0\pi$. Пусть $n=\dim\tau$~---


размерность базы расслоения $\pi$ и


$r=\dim\nu$~--- ранг $\pi$. По предложению \ref{proppre}


\begin{multline*}


\dim\g^l(x_k)=n\bigg(\tbinom{l+r-1}{r-1}+\sum_{1\le


i<k}\tbinom{i+n-1}{n-1}\tbinom{l+r-2}{r-1}+


\sum_{i=k}^{k+l-1}\tbinom{i+n-1}{n-1}\tbinom{l+k+r-i-2}{r-1}\bigg)+\\


+r\bigg(\sum_{0\le i<k}\tbinom{i+n-1}{n-1}\tbinom{l+r-1}{r-1}+


\sum_{i=k}^{l+k}\tbinom{i+n-1}{n-1}\tbinom{l+k+r-i-1}{r-1}\bigg).


\end{multline*}


Используя асимптотику $\binom{l+s}s\sim\frac{l^s}{s!}$ при


$l\to\infty$, получаем


\begin{equation}


\label{fo1}


\dim\g^l\sim (n+r)\frac{l^{n+r-1}}{(n+r-1)!}.


\end{equation}





\begin{prop}


Единственными уравнениями, удовлетворяющими условию $(\ref{reform-trans})$,


являются уравнения размерности $d=n+r$ и порядка $k=1$. Более того,


возможны только следующие случаи\/${:}$ $1)$~$n=1$,


$2)$~$r=1$, $3)$~$n=r=2$.


\end{prop}





\begin{pf}


Напомним, что размерность пространства струй $J^k\pi$ равна


$d_k=n+r\binom{n+k}k$. Пусть $d=\dim\E$, тогда


$$


\dim h^l_{x_k}=\dim(S^lT^*_{x_k}\E\ot\vg_{x_k})=(d_k-d)\cdot


\tbinom{d+l-1}l\sim(d_k-d)\cdot\dfrac{l^{d-1}}{(d-1)!}.


$$


Сравнивая эту асимптотику с (\ref{fo1}) и (\ref{reform-trans}), получаем


$d\le n+r$, а т.\,к. $\pi_{k,0}:\E\to J^0\pi$ сюръективно, то на самом


деле выполняется равенство. Кроме того, справедливо неравенство


$n+r\ge r[\binom{n+k}k-1]$, которое при $k>1$ имеет только одно


решение $k=2$, $n=r=1$. Так размерность подмногообразия


$\E\subset J^2(1,1)$ равна


двум, оно задается одним уравнением первого порядка и одним уравнением


второго порядка, которое является продолжением первого


уравнения. Поэтому, на самом деле, $\E$ есть уравнение первого порядка.





Таким образом, $k=1$, следовательно,


$d_k-d=nr$. Решения дополнительного неравенства $n+r\ge nr$ имеют вид


1)~$n=1$; 2)~$r=1$; 3)~$n=r=2$. Для этих значений неравенство


(\ref{reform-trans}) выполняется для всех $l$ точно (а не только


асимптотически).


\end{pf}








Теперь рассмотрим три полученных случая.





1. $n=k=1$. Подмногообразие $\E\subset J^1(1,r)$


коразмерности $r$ есть определенная система обыкновенных


дифференциальных уравнений. Согласно известной теореме существования и


единственности все такие системы локально эквивалентны, т.\,е.


получаем трансверсальность.





2. $r=k=1$. Здесь $\E\subset J^1(n,1)$ коразмерности $n$


диффеоморфно проектируется с помощью $\pi_0$ на $J^0\pi$. Через


каждую точку $x_0=\pi_0(x_1)$ проходит $n$-плоскость


$L(x_1)$. Полученное распределение ранга $n$ на многообразии


$\E$ размерности $n+1$, вообще говоря, не интегрируемо (что


соответствует неинтегрируемости $\E$) и является либо контактным,


либо четно-контактным. В обоих случаях получаем трансверсальность и,


следовательно, локальную эквивалентность всех таких уравнений.





3. $k=1$, $n=r=2$. Используя комплексные координаты, получаем


обыкновенные (не обязательно голоморфные) дифференциальные уравнения


$w'_z=\psi(z,\ov z,w,\ov w)$. Вновь получаем трансверсальность и


локальную эквивалентность.





\begin{note}\label{rkrk}


Так как $\pi_{1,0}:\E\to J^0\pi$ является диффеоморфизмом, получаем


распределение $L(\pi_{1,0}^{-1}(\cdot))$ на $J^0\pi$. Поэтому


описанные случаи соответствуют известным распределениям, не имеющим


модулей: 1)~одномерное распределение; 2)~(четно-)контактное


распределение коранга~1; 3)~распределение Энгеля с


$\rank=\op{corank}=2$.


\end{note}





Заметим, что если ограничиться эквивалентностями


индуцированными морфизмами расслоения $\pi$ (с порождающими полями


вида $X=\sum a^i(x)\p_{x^i}+b^j(x,u)\p_{u^j}$), то соответствующие размерности


имеют асимптотику: $\dim\g^l\sim


r\cdot\frac{l^{n+r-1}}{(n+r-1)!}$.


Поэтому неравенство (\ref{reform-trans}) может выполняться в


единственном случае $n=1$, т.\,е. для обыкновенных


дифференциальных уравнений первого порядка. На самом деле, в этом


случае мы имеем трансверсальность также относительно меньшей группы.


\medskip





5.3. {\it Псевдогруппа Ли контактных преобразований}.


В случае $\rank\pi=1$ вычисления аналогичны и из предложения


\ref{proppro} следует асимптотика


\begin{equation}


\label{fo4}


\dim\g^l\sim \dfrac{l^{2n}}{(2n)!}.


\end{equation}





\begin{prop}


Единственными уравнениями, удовлетворяющими условию


$(\ref{reform-trans})$, являются те, которые имеют размерность


$n+1\le d\le 2n+1$, и либо


$\pi_1(\E)\subset J^1\pi$ собственное, либо $n=1$, $k=2$.


\end{prop}





\begin{pf}


Первое утверждение получается, если подставить  (\ref{fo4}) в


(\ref{reform-trans}) и учесть, что $\pi_{k,0}|_\E$ сюръективно.


Тем не менее, если $d=2n+1$, то как в случае $\pi_{k,1}(\E)=J^1\pi$,


получаем дополнительное равенство $d_k-d=1$, имеющее единственное решение


$n=1$, $k=2$.


\end{pf}








Обсудим два полученных случая.





1. $k=1$ и $\E\subset J^1(n,1)$. Как и в п.\,3.5, видим,


что дифференциальное уравнение в частных производных


$\E\subset J^1\pi$, $\dim\E=d$, трансверсально относительно


псевдогруппы Ли контактных преобразований в $x_1\in\E$ тогда и только


тогда, когда не существует интегральных подмногообразий размерности


большей чем $\big[\frac{d-1}2\big]$. Заметим, что распределение


$\Pi\cap T\E$, индуцированное на $\E$, всегда обладает интегральными


подмногообразиями $L$ размерности $\big[\frac{d-1}2\big]$.


Если $\pi_1:L\to L_0$~--- диффеоморфизм, подмногообразие имеет вид


$j_1s(L_0)$ для некоторого сечения $s$ расслоения $\pi$. Поэтому


трансверсальность $\E$ означает, что не существует ``частных решений''


$s:L_0\to J^0\pi$,


$j_1s(L_0)\subset\E$, размерности большей, чем минимально возможная.





2. $k>1$. Если $\E_1=\pi_{k,1}(\E)\subset J^1\pi$


собственное, то либо не существует продолжения $\E_1^{(1)}$, либо


уравнение $\E_1$, а значит, и $\E$ не трансверсальны. Поэтому рассмотрим


$\E_1=J^1\pi$, тогда $\pi_{2,1}:\E_2\to J^1\pi$~--- диффеоморфизм.


Как показано в предыдущем предложении, тогда $n=1$. В этом случае


имеем $l$-трансверсальность для всех $l$.





На самом деле, это~--- известный результат Софуса Ли: все скалярные


обыкновенные дифференциальные уравнения контактно эквивалентны.





Таким образом, справедлива





\begin{thm}


Единственными трансверсальными $($и эквивалентными$)$ уравнениями $\E\subset


J^k\pi$ относительно псевдогруппы преобразований Ли являются следующие\/${:}$


\begin{itemize}


\item[1.] $u'_x=F(x,u)$, $x\in\R$, $u\in\R^n;$


\item[2.] $u'_{x^i}=\vp_i(x,u)$, $i=1,\dotsc,n$, $x\in\R^n$, $u\in\R;$


\item[3.] $w'_z=\vp(z,\ov z,w,\ov w)$, $w'_{\ov z}=\psi(z,\ov z,w,\ov


w)$, $z,w\in\CC;$


\item[4.]


$u'_{x^i}=\vp_i(x,u,u'_{x^{s+1}},\dotsc,u'_{x^n})$, $1\le i\le


s<n$, $x\in\R^n$, $u\in\R;$


\item[5.]


$u''_{xx}=F(x,u,u'_x)$, $x\in\R$, $u\in\R$.


\end{itemize}


\end{thm}





Согласно замечанию \ref{rkrk}, первые три случая соответствуют


открытым орбитам в пространстве ростков распределений.


Следующий случай описывается как подмногообразие контактного многообразия.


Последний эквивалентен слоению Лежандра на контактном многообразии $J^1(1,1)$.


На самом деле, локально все лежандровы слоения на контактном


многообразии эквивалентны, но только при $n=1$ (согласно теореме


С.\,Ли) соответствующее уравнение $\E$ общего положения. В противном


случае, необходимо наложить на $\E$ дополнительное условие


интегрируемости. Для таких уравнений справедливо





\begin{prop}


Псевдогруппа преобразований Ли транзитивно действует на классе


$\frak{N}$ уравнений $\E\subset J^{1+\varepsilon}\pi$ $($где


$\varepsilon=\max(0,2-r)$ как и раньше$)$, которые интегрируемы и


диффеоморфно проектируются с помощью


$\pi_{1+\varepsilon,\varepsilon}:\E@>\sim>>\to J^\varepsilon\pi$.


\end{prop}





\begin{pf}


Действительно, при $r=\rank\pi=1$ получаем лежандровы слоения


контактного многообразия $J^1\pi$, а при $r>1$~--- просто слоение на


$J^0\pi$.


\end{pf}








\section{Заключение}





Мы описали препятствия к трансверсальности, а значит, и к


эквивалентности, относительно действия псевдогруппы $G$. Оказывается,


что эти препятствия имеют когомологическую природу. Поясним это для


пространств ${\frak O}^l$, возникающих в критерии трансверсальности


${\frak O}^l_a=0$ в предложении \ref{thm4}.





Обозначим через $H^{l-j,j}({\frak h})$ группу когомологий комплекса


$$


%%\begin{equation}\label{hhh2}


0 \to{\frak h}^l_a \to {\frak h}^{l-1}_a\ot\tf^*_a @>\d>>


{\frak h}^{l-2}_a\ot\La^2\tf^*_a @>\d>> \dotsb\to {\frak


h}^{l-j}_a\ot\La^j\tf^*_a@>\d>> \dotsb


$$


Заметим, что $\d$-дифференциал действует вдоль $\tf$, а не вдоль $TM$,


поэтому группы $H^{l,0}({\frak h})={\frak h}_l\cap


S^l(\Ann\tf)\ot TM$ могут не быть нулевыми. Тем не менее, они


обращаются в нуль, если подпространство $\tf_a$ (слабо) не


характеристическое.





Если псевдогруппа $G$ порядка $k$ является $2$-ацикличной, то


первая группа когомологий вышеуказанного комплекса тривиальна, поэтому


${\frak h}^l_a$ продолжает ${\frak h}^{l-1}_a$ вдоль $\tf$ при


$l>k$. Также предположим, что $l>m$, где $m$~--- порядок уравнения на


подмногообразии ${\frak N}$, тогда $h_a^l=(h_a^{l-


1})^{(1)}$, $\g_a^l=(\g_a^{l-1})^{(1)}$ при $l>k$.





При использовании стандартных методов гомологической алгебры имеет место





\begin{thm}


Пусть псевдогруппа $G$ \ $3$-ациклична. Тогда, если


${\frak O}_a^{l-1}=0$, то ${\frak O}_a^l=H^{l-2,2}


(\frak{h}_a)$. В частности, если действие $G$ трансверсльно


порядка $l$ и $\frak{h}$ \ $2$-ациклична, то действие формально


трансверсально. 


\end{thm}





Так обстоит дело, например, для комплексных, симплектических и некоторых


других псевдогрупп, для которых трансверсальность была получена


другими методами.








\begin{thebibliography}{99}








\bibitem{Lie} %1


Lie S.


{\em Theorie der Transformationsgruppen}. -- Zweiter Abschnitt, unter


Mitwirkung von Prof. Dr.~Friederich Engel, Leipzig: Teubner, 1890.





\bibitem{C} %2


Cartan E. {\it Sur la structure des groupes infinis de


transformations}. -- in: {Oeuvres completes}, 


Paris: Gautier-Villars, 1953. -- V.\,II. -- \No\,2, -- P.\,571--714.





\bibitem{E} %3


Ehresmann C. {\it Introduction \`a la th\'eorie des structures


infinit\'esimales et des pseudogroupes de Lie} // Colloq. internat.


Centre nat. rech. scient. 52. --


Strasbourg, 1953. -- \No\,11. -- P.\,97--110.





\bibitem{Lib} %4


Libermann P. {\it Pseudogroupes infinit\'esimaux de Lie} //


C.\;R. Acad. Sci. Paris. -- 1958. -- V.\,246. -- P.\,531--534.





\bibitem{SS} %5


Singer I.M., Sternberg S. {\it On the infinite groups of Lie


and Cartan} // J. d'Analyse Math. -- 1965. -- V.\,15. -- P.\,1--114.





\bibitem{M} %6


Malgrange B. {\it Equations de Lie}. I // J. Diff.


Geom. -- 1972. -- V.\,6. -- P.\,503--522.





\bibitem{M1} %6a


Malgrange B. {\it Equations de Lie}. II // J. Diff.


Geom. -- 1972. -- V.\,7. -- P.\,117--141.








\bibitem{KLV} %7


Виноградов А.М., Красильщик И.С., Лычагин В.В. {\em Введение в геометрию


нелинейных


дифференциальных уравнений}. -- М.: Наука, 1986. -- 335 c.





\bibitem{S} %8


Spencer D. {\it Deformation of structures on manifolds defined


by transitive pseudogroups} // Ann. of Math. -- 1962. -- V.\,76. --


\No\,2. -- P.\,306--445.





\bibitem{BM} %9


Buttin C., Molino P. {\it Theoreme general d'equivalence pour


les pseudogroupes de Lie plats transitifs} // J. Diff.


Geom. -- 1974. -- V.\,9. -- P.\,347--354.





\bibitem{P} %10


Pollack A.S. {\it The integrability problem for pseudogroup


structures} // J. Diff. Geom. -- 1974. -- V.\,9. -- P.\,355--390.





\bibitem{GS} %11


Guillemin P., Sternberg S. {\it Deformation theory of pseudogroup


structures} // Mem. AMS. -- 1966. -- V.\,64. -- P.\,1--80.





\bibitem{T} %12


Tanaka N.


{\it On differential systems, graded Lie algebras and pseudogroups} //


J. Math. Kyoto Univ. -- 1970. -- V.\,10. -- P.\,1--82.





\bibitem{Ly} %13


Lychagin V. {\it Homogeneous geometric structures and


homogeneous differential equations} // AMS. Transl. Ser.


The interplay between differential geometry and differential


equations. Ser.~2. -- 1995. -- V.\,167. -- P.\,143--164.





\bibitem{Gu} %14


Guillemin P.


{\it The integrability problem for $G$-structures} //


Trans. AMS. -- 1965. -- V.\,116. -- P.\,544--560.





\bibitem{Go} %15


Goldschmidt H. {\it Integrability criteria for systems of


nonlinear partial differential equations} // J. Diff. Geom. -- 1967. --


V.\,1. -- P.\,269--307.





\bibitem{Ko} %16


Кобаяси Ш. {\em Группы преобразований в дифференциальной геометрии}. --


М.: Наука, 1986. -- 224 c.





\bibitem{CM} %17


Chern S.S., Moser J. {\it Real hypersurfaces in complex


manifolds} // Acta Math. -- 1974. -- V.\,133. -- P.\,219--271.





\bibitem{Kr1} %18


Кругликов Б.С. {\em Тензоры Ниенхейса и препятствия к построению $1$


псевдоголоморфных отображений} // Матем. заметки. -- 1998. -- T.\,63.


-- Bып.\,4. -- C.\,541--561.





\bibitem{Kr2} %19


Kruglikov B.S. {\it Non-existence of higher-dimensional


pseudoholomorphic submanifolds} // Manuscripta Mathematica. -- 2003.


-- V.\,111. -- P.\,51--69.





\bibitem{AG} %20


Арнольд В.И., Гивенталь А.Б. {\em Симплектическая геометрия} //


Итоги науки и техн. Современ. пробл. матем.


-- М.: ВИНИТИ, 1985. -- Т.\,4. -- C.\,5--139.








\end{thebibliography}





\vskip 2cm





\post{\it Университет Тромсо}{\it Поступила}


\post{\it $($Норвегия$)$}{02.03.2004}





\label{end}


\end{document}








Mat-Stat. Dept., University of Troms\o, Norway








\bibitem{KL}


Kruglikov B.S., Lychagin V.V.


{\it On equivalence of differential equations} //


Acta et Comment. Univ. Tartuensis Math. -- 1999. -- V.\,3. -- P.\,7--29.





